\section{The \bench Benchmark}
\label{sec:bench}

\subsection{Overview}
\label{sec:bench_overview}
A dynamic benchmark must charge each policy for the time it takes to compute. When the
simulation advances only after the policy has returned its actions, as in the default
evaluators of DOMINO \cite{domino} and PhysMani-Bench \cite{physmani}, a slow policy is
scored as if it were instant: the evaluation measures what a policy predicts but not
whether the prediction arrives in time, which is what decides success on a real robot.

To evaluate whether a policy can act on a moving scene in time, we introduce \bench, a
benchmark of 14 dynamic manipulation tasks in which the world never pauses.
Observation, scene motion, and action execution share one physical clock, so the scene
keeps moving for as long as each policy computes. Fig.~\ref{fig:bench_overview}(a) shows what this changes: during a 200\,ms inference, a
rolling ball waits at its observed position when the world is paused, but in \bench it has
rolled 18\,cm, so an action aimed at the observed position misses. \bench covers five task
families (Fig.~\ref{fig:bench_overview}(b)): \emph{Grasp}, grasping a target under
harmonic, conveyor, translating, rotating, or occluded motion; \emph{Place}, placing a
held object into a translating or rotating receptacle; \emph{Pass}, passing a wand
along a rotating wire; \emph{Capture}, capturing a rolling or sliding object; and
\emph{Impact}, grasping a part after an impact. Fig.~\ref{fig:bench_overview}(c) shows
all 14 tasks. Together they vary what decides success: the target's speed, from 0.03
to 1.14\,m/s, its rotation, the clearance of obstacles, and the motion imparted by
drivers such as conveyors, ramps, and strikers.

\subsection{Construction}
\label{sec:bench_construction}
\bench is built in MuJoCo and robosuite \cite{mujoco,robosuite} with a Franka Panda arm
controlled at \cfgControlRate{}. Each scene samples the target's initial state and
motion parameters within task-specific ranges. A task-specific expert with privileged
simulator state provides \cfgDemos{} demonstrations per task, 11{,}200 in total. All
results are reported on 1{,}400 held-out test scenes, 100 per task, shared by all
methods; configurations are selected on separate development scenes.

\subsection{Evaluation Protocol}
\label{sec:bench_protocol}
\paragraph{Native latency.}
During each query, physics advances by the method's own inference latency, drawn from
its latency distribution measured at batch size 1 on one \cfgGPU{}. All methods share
one executor, which hands over each chunk at the first control-period boundary after
inference completes and executes it from the command indexed to that boundary.

\paragraph{Settings and metrics.}
We evaluate three settings: \emph{Native}, at each method's own latency; \emph{+150\,ms},
with 150\,ms added to every method; and \emph{Paused}, with the world paused during inference. We report the task-balanced success rate (SR) and
the p95 time-to-action, the time from capturing an observation until its first command is available. The full
protocol is in the supplementary material.
